% 07-implementacion.tex — Fuente: docs/informe-final/07-implementacion.md

\chapter{Implementación}
\label{cap:implementacion}

No se transcriben archivos completos; se remite al repositorio para el
detalle. Este capítulo ilustra con listados pequeños y significativos las
decisiones críticas de código.

\section{Manejo de errores globales (RFC 7807 Problem Details)}

Todos los errores de la API se serializan como \texttt{ProblemDetail}
(RFC~7807), en vez de estructuras JSON ad-hoc distintas por controlador, como se puede observar en el Listado~\ref{lst:manejador-excepciones}.

\begin{lstlisting}[language=Java,caption={Manejador global de excepciones (RFC 7807 ProblemDetail). Fuente: \texttt{artisync/\allowbreak Backend/\allowbreak src/\allowbreak main/\allowbreak java/\allowbreak .../\allowbreak exception/\allowbreak ManejadorGlobalExcepciones.java}},label=lst:manejador-excepciones]
@Slf4j
@RestControllerAdvice
public class ManejadorGlobalExcepciones {

    private static final String BASE_TIPO = "https://artisync.dev/errors/";

    @ExceptionHandler(AccessDeniedException.class)
    public ResponseEntity<ProblemDetail> manejarAccesoDenegado(AccessDeniedException ex) {
        ProblemDetail problema = ProblemDetail.forStatus(HttpStatus.FORBIDDEN);
        problema.setType(URI.create(BASE_TIPO + "acceso-denegado"));
        problema.setTitle("Acceso denegado");
        // ... (continua con los demas @ExceptionHandler: validacion, JWT, 404, etc.)
    }
}
\end{lstlisting}

Esta decisión centraliza el formato de error para las 1230 pruebas del
backend y para el interceptor HTTP del frontend Angular, evitando el
manejo de errores ad-hoc por controlador.

\section{Esquema de caching}

El listado de catálogo (endpoint de mayor frecuencia de lectura,
\texttt{REQ-NF-004}) usa Spring Cache sobre Redis (\texttt{ADR-004}), con
invalidación automática (\texttt{@CacheEvict}) al crear/editar/eliminar
un servicio, tal como se muestra en el Listado~\ref{lst:cache-catalogo}:

\begin{lstlisting}[language=Java,caption={Caché del listado de catálogo (Spring Cache + Redis). Fuente: \texttt{artisync/\allowbreak Backend/\allowbreak src/\allowbreak main/\allowbreak java/\allowbreak .../\allowbreak service/\allowbreak catalogo/\allowbreak impl/\allowbreak ServicioCatalogoServicioImpl.java}},label=lst:cache-catalogo]
@Override
@Transactional(readOnly = true)
@Cacheable(cacheNames = "catalogo")
public Page<RespuestaServicioResumido> buscarCatalogoServicios(
        Long categoriaId, Long subcategoriaId, BigDecimal precioMin, BigDecimal precioMax,
        List<Long> etiquetaIds, String textoBusqueda, String sortParam, int page, int size) {
    // ... construye la Specification dinamica y ejecuta la consulta paginada
}
\end{lstlisting}

La clave de caché de Spring se compone de todos los parámetros del método
--- una limitación de diseño relevante para la interpretación de los
resultados de rendimiento del Capítulo~\ref{cap:evaluacion-resultados}
(ver \texttt{docs/mediciones/perf/REPORTE-PERF.md}, advertencia
metodológica sobre el escenario ``frío'').

\section{Mecanismo de seguridad: JWT + rate limiting}

La autenticación es \emph{stateless} (JWT~\cite{JWT2015}, \texttt{ADR-002});
las rutas de autenticación aplican rate limiting por IP real (resuelta
vía \texttt{server.forward-headers-strategy=native}, no
\texttt{getRemoteAddr()} directo) y por cuenta, con respuesta \texttt{429}
en formato \texttt{ProblemDetail} --- ver \texttt{AuthRateLimitFilter} e
\texttt{IntentosAutenticacionService}, corrección de la observación
\texttt{OBS-AUTO-05} documentada en \texttt{docs/observaciones/OBSERVACIONES.md}.
El diagrama de la Figura~\ref{fig:secuencia-jwt} ilustra la secuencia.

\begin{figure}[H]
\centering
\includegraphics[width=0.85\textwidth]{secuencia_login_jwt.png}
\caption{Diagrama de secuencia del flujo de autenticación JWT. Fuente:
\texttt{docs/diagramas/secuencia\_login\_jwt.png}.}
\label{fig:secuencia-jwt}
\end{figure}

\section{Estrategia híbrida de acceso a datos: procedimiento almacenado representativo}
\label{sec:acceso-datos-hibrido}

Esta subsección ilustra el contrato completo entre un procedimiento
versionado en \texttt{db/procs/} y el método de repositorio Spring Data
que lo invoca, siguiendo la especificación JPA~2.1 (\texttt{ADR-006},
Capítulo~\ref{cap:diseno-arquitectura}~\S6.2).
El Listado~\ref{lst:sp-registrar-decision} presenta la definición SQL, mientras que el Listado~\ref{lst:repo-registrar-decision} detalla su invocación formal con la anotación \texttt{@Procedure}.

\begin{lstlisting}[language=SQL,caption={\texttt{sp\_registrar\_decision\_verificacion} (categoría: validaciones cruzadas). Fuente: \texttt{artisync/\allowbreak Backend/\allowbreak src/\allowbreak main/\allowbreak resources/\allowbreak db/\allowbreak migration/\allowbreak V7\_\_verificacion\_asistida\_ia.sql}},label=lst:sp-registrar-decision]
CREATE OR REPLACE PROCEDURE sp_registrar_decision_verificacion(
    p_id_certificado BIGINT,
    p_id_estado BIGINT,
    p_id_moderador BIGINT,
    p_nota TEXT
)
LANGUAGE plpgsql
AS $$
BEGIN
    UPDATE certificados_ia 
    SET id_estado = p_id_estado, id_moderador = p_id_moderador, nota = p_nota, fecha_analisis = NOW()
    WHERE id_certificado = p_id_certificado;
END;
$$;
\end{lstlisting}

\begin{lstlisting}[language=Java,caption={Método de repositorio que invoca \texttt{sp\_registrar\_decision\_verificacion}. Fuente: \texttt{artisync/\allowbreak Backend/\allowbreak src/\allowbreak main/\allowbreak java/\allowbreak .../\allowbreak repository/\allowbreak profile/\allowbreak AiCertificateRepository.java}},label=lst:repo-registrar-decision]
    @Procedure(procedureName = "sp_registrar_decision_verificacion")
    void registrarDecision(
            @Param("p_id_certificado") Long idCertificado,
            @Param("p_id_estado") Long idEstado,
            @Param("p_id_moderador") Long idModerador,
            @Param("p_nota") String nota);
\end{lstlisting}

\subsection{Conformidad lograda}

La invocación utiliza los mecanismos formales de la
especificación JPA~2.1 mediante la anotación \texttt{@Procedure} sobre el método del
repositorio Spring Data, asegurando una conexión segura y \textbf{sin
concatenación de cadenas} (evitando así la inyección SQL, conforme al
análisis de Gould~\cite{GOULD2004} y Halfond~\cite{HALFOND2005}). Este
esquema resuelve plenamente los requerimientos de seguridad y diseño arquitectónico (ver Sección~\ref{sec:acceso-datos-hibrido}).

\subsection{Resolución aplicada: rutinas con retorno escalar}

\texttt{@Procedure} se aplica de forma sistemática cuando el método Java es
\texttt{void} y la rutina no declara ningún parámetro \texttt{OUT}/\texttt{INOUT}: además de
\texttt{sp\_registrar\_decision\_verificacion} (Listado~\ref{lst:repo-registrar-decision}), es el
caso de \texttt{sp\_restablecer\_contrasena} y \texttt{sp\_cambiar\_contrasena}, migradas y
verificadas contra el stack real (revisión técnica 2026-09-05).

Para las rutinas que necesitan devolver un valor escalar --- por ejemplo
\texttt{fn\_registrar\_usuario}, \texttt{fn\_resolver\_estado\_login} o
\texttt{fn\_permisos\_efectivos\_usuario} (esta última se ejecuta en cada inicio de sesión) --- se
investigó en tres rondas independientes (2026-09-05, 2026-09-14 y una verificación adicional la
misma noche) si podían invocarse con \texttt{@Procedure}/\texttt{@NamedStoredProcedureQuery}. Las
tres coincidieron en el mismo resultado: en cuanto un método declara un tipo de retorno distinto de
\texttt{void}, Hibernate~7.4.1 serializa \emph{todos} los parámetros --- incluido el propio valor de
retorno --- con la sintaxis de argumento nombrado de PostgreSQL dentro del \emph{JDBC escape}
\texttt{\{call ...\}}, que PostgreSQL no puede interpretar ahí
(\texttt{ERROR: syntax error at or near "=>"}). Se confirmó que ni los parámetros puramente
posicionales evitan el defecto (Spring Boot compila con \texttt{-parameters} y Hibernate obtiene el
nombre por reflexión de todos modos), ni una versión de parche posterior de Hibernate
(\texttt{7.4.5.Final}, probada explícitamente contra Postgres real y revertida sin dejar rastro) lo
corrige. La única vía que evita el defecto --- la API nativa imperativa de Hibernate
(\texttt{Session.createStoredProcedureCall} + \texttt{markAsFunctionCall}) --- es estructuralmente
inalcanzable desde cualquier anotación JPA declarativa: exigiría abandonar
\texttt{@Procedure}/\texttt{@NamedStoredProcedureQuery} por completo e implementar cada método a
mano. Un commit del 4 de septiembre de 2026 que migró estas rutinas a \texttt{@Procedure} sin
conocer aún esta restricción rompió el inicio de sesión en producción; se revirtió y motivó esta
investigación. El detalle completo, con la orden y la salida de cada prueba contra Postgres real,
está en \texttt{docs/basedatos/CATALOGO-SP.md} (\S14).

\textbf{Resolución (14-sep-2026, cierre del Punto~P6 de la guía del examen suspenso):} la guía
exige, literalmente, \emph{cero apariciones de \texttt{nativeQuery=true} para invocar
procedimientos almacenados} --- no exige la anotación \texttt{@Procedure} en sí. Puesto que ninguna
vía dentro de JPA/Hibernate evita el defecto sin renunciar a una anotación declarativa, la solución
fue salir de JPA/Hibernate por completo para estas 23 rutinas: se sustituyó
\texttt{@Query(nativeQuery = true)} por \texttt{NamedParameterJdbcTemplate}
(\texttt{org.springframework.jdbc.core.namedparam}), que construye la sentencia como texto
(\texttt{SELECT fn\_x(:parametro)}) y la ejecuta por JDBC puro, sin instanciar nunca un
\texttt{Session}/\texttt{EntityManager} de Hibernate --- y por tanto sin pasar por las clases donde
se origina el defecto (\texttt{PostgreSQLCallableStatementSupport},
\texttt{NamedCallableQueryMementoImpl}). El patrón aplicado es consistente con código ya existente
en el proyecto (\texttt{ContractRepositoryCustom}/\texttt{ContractRepositoryImpl}, que ya resolvía
consultas dinámicas complejas por fuera de \texttt{@Query}): cada repositorio afectado gana una
interfaz \texttt{XxxRepositoryCustom} y una clase \texttt{XxxRepositoryImpl} que la implementa
inyectando \texttt{NamedParameterJdbcTemplate}; para quien invoca el repositorio no cambia nada,
Spring Data fusiona ambas interfaces por convención de nombre.

Se verificó que este cambio no reproduce el incidente de producción del 4 de septiembre de 2026 que
motivó originalmente la excepción: suite completa de pruebas (1448/1448), 31 pruebas de integración
contra Postgres real, y login, registro, recuperación de contraseña y gestión de roles/países
probados a mano contra el backend con Postgres real, incluida la captura del \texttt{Set-Cookie}
real del login.

De paso se revisaron también los dos \texttt{PROCEDURE} de purga
(\texttt{sp\_purgar\_notificaciones}, \texttt{sp\_purgar\_datos\_seguridad}, invocados hoy con
\texttt{JdbcTemplate.update("CALL ...")} desde \texttt{NotificationPurgeScheduler} y
\texttt{SecurityPurgeScheduler}): son \texttt{void}, solo parámetros \texttt{IN}, exactamente el
patrón que sí funciona con \texttt{@Procedure}, pero ambos ejecutan \texttt{COMMIT} explícito por
lote dentro del propio \texttt{PROCEDURE} (para no sostener una transacción larga que bloquee
\texttt{VACUUM}), lo que exige que la conexión no tenga ya una transacción abierta por Spring
(\texttt{Propagation.NOT\_SUPPORTED} está declarado como obligatorio en ambos \textit{schedulers},
no cosmético) --- Postgres falla con \texttt{2D000 invalid\_transaction\_termination} si el
\texttt{CALL} ocurre dentro de una transacción ya iniciada. Se mantienen en \texttt{JdbcTemplate}
directo --- ya usaban este mismo mecanismo antes de la migración de las otras 23 rutinas --- no por
descuido sino por esta restricción de manejo de transacciones, independiente del defecto de
Hibernate.

De las 28 rutinas activas (26 de \texttt{db/procs/} más 2 de verificación asistida por IA), el
reparto exacto por mecanismo de invocación --- verificable de forma reproducible con
\texttt{python scripts/auditoria-rubrica.py p6}, que descuenta las menciones dentro de comentarios
Javadoc --- es: \textbf{3} vía \texttt{@Procedure} (método \texttt{void}, sin parámetros
\texttt{OUT}/\texttt{INOUT}), \textbf{0} vía \texttt{@NamedStoredProcedureQuery}, \textbf{0} vía
\texttt{@Query(nativeQuery = true)} y \textbf{25} vía \texttt{NamedParameterJdbcTemplate} (las 23
migradas en esta resolución más las 2 rutinas de purga, que ya usaban este mecanismo desde antes por
la restricción de transacciones descrita arriba). El mecanismo de acceso a datos, en conjunto, sigue
siendo \emph{híbrido y documentado por diseño} (ADR-006), no accidental ni pendiente de resolución.
